若未使用 AI, 需要在参考文献前使用命令~\verb|\CumcmAINotUsed| 加入如下段落:

\CumcmAINotUsed

若使用了 AI, 需要使用命令 \verb|\CumcmAIUsed{<用途>}| 加入如下段落:
\CumcmAIUsed{[简要用途, 如语言润色, 代码调试等]}

支撑材料格式如下.

\section*{AI工具使用详情}

\subsection*{工具信息}
\noindent
\begin{tabularx}{\textwidth}{@{}XX@{}}
  \toprule
  项目 & 填写内容 \\
  \midrule
  工具名称 & \textit{填写实际使用的产品或工具名称} \\
  版本或型号 & \textit{填写实际版本、模型名称或可识别的型号} \\
  使用日期 & \textit{YYYY-MM-DD 至 YYYY-MM-DD} \\
  \bottomrule
\end{tabularx}

\subsection*{具体使用目的和环节}

逐项说明 AI 用于语言润色、代码调试、资料整理或其他环节；不得笼统写“辅助论文”。明确其输入、输出及对应论文位置。

\subsection*{主要提示方式与使用过程}

说明主要提示策略与迭代过程，可附具有代表性的交互片段。不要包含参赛者姓名、学校或赛区信息。

\subsection*{采纳、人工修改与核验}

除语言润色外，逐项记录 AI 输出是否采纳、由谁进行何种人工修改、采用何种数据或实验复核，以及最终结论是否发生变化。

\subsection*{人工核验清单}

\begin{enumerate}[label=\arabic*)]
  \item 核对公式推导、单位、边界条件和符号一致性；
  \item 在独立环境中运行全部代码并复现论文表图；
  \item 核验引用、数据来源和事实陈述；
  \item 确认核心建模与分析由参赛队主导；
  \item 确认论文声明与本详情文件一致。
\end{enumerate}
